\section{Experimental data and observers}\label{sec:experiment}
We use $\TV(P,Q)=\sup_B|P(B)-Q(B)|=\frac12\|P-Q\|_1$. Fix a finite horizon $n\ge1$. A call takes a hidden physical state $x$, a legal action $a$, and produces a new hidden state, a report, and a resource increment through a positive normalized kernel. Positive means nonnegative and normalized; deterministic reports and zero likelihoods are permitted. Preparation is fixed. An unknown parameter, when present, is not an argument available to the observer. This section first defines a single prepared law. Adjoining an unknown parameter to the hidden state gives a Bayes experiment, not automatically a frequentist minimax theorem. Section~\ref{sec:uniform} supplies the additional common-programme conditions and quantifiers for the latter.

An executable future test is a legal continuation with a bounded scored outcome. At a fixed phase, histories are equivalent if all such tests have the same conditional expectation and the same legal interface. Appendix~\ref{sec:quotient} gives sufficient raw-kernel hypotheses for a compact metrizable realization of this quotient. In particular it establishes measurability, actual mixture realizability, and minimality. We now state precisely the properties used by the main theorem.

For each $t=0,\ldots,n$, let $K_t$ be a compact metrizable convex subset of a Hausdorff locally convex space, embedded by a countable separating family of continuous affine future-test expectations. Write $\aff(K_t)$ for continuous affine functions, including constants. Probability measures on $K_t$ have barycenters in $K_t$, denoted $\bar\nu$, and $f(\bar\nu)=\int f\,d\nu$ for $f\in\aff(K_t)$. No finite dimension is assumed. Let $\mathcal A$ be finite and let $\mathcal A_t\subseteq\mathcal A$ be a fixed nonempty legal action set at phase $t<n$. Let $\mathcal Y_a$ be standard Borel. These sets are part of the physical interface, not a phase variable readable by a stationary machine. Kernels for illegal actions may be extended arbitrarily and never enter a positive-mass term. For $t<n$, the raw instruments induce jointly Borel maps
\begin{equation}\label{eq:instrument}
 J_t(q,a)(dy),\qquad F_t(q,a,y)\in K_{t+1},\qquad
 \Lambda_t(q,a)=\law_{Y\sim J_t(q,a)}F_t(q,a,Y).
\end{equation}
Values of $F_t$ on zero-likelihood reports are irrelevant and are assigned measurably. The instrument property means that, for every probability $\nu$ on $K_t$, Borel report set $B$, and $f\in\aff(K_{t+1})$,
\begin{align}\label{eq:affine_instrument}
 J_t(\bar\nu,a)(B)&=\int J_t(q,a)(B)\nu(dq),\\
 \int_B f(F_t(\bar\nu,a,y))J_t(\bar\nu,a)(dy)
 &=\int\!\int_B f(F_t(q,a,y))J_t(q,a)(dy)\nu(dq).\notag
\end{align}
The identities are required on all mixture states in the carrier. They express physical mixing before an unnormalized instrument, not linearity of the normalized Bayes update. Assume $q\mapsto\Lambda_t(q,a)$ is weakly continuous. This ensures continuous Bellman values; the feasibility parts below only need the displayed measurable instruments and conditional barycenters. Total-variation report continuity is an additional quantitative property verified in the regular realizations, not a premise of convex order.

The terminal decisions form a compact metric space $U$. The scored loss has conditional expectation $\ell(q,u)\in[0,H]$, jointly continuous and affine in $q$ for each $u$. Let
\begin{equation}\label{eq:bellman}
 g(q)=\min_{u\in U}\ell(q,u),\quad V_n=g,\quad
 V_t(q)=\min_{a\in\mathcal A_t}\int V_{t+1}(w)\Lambda_t(q,a)(dw),\quad B_n^*=V_0(q_0).
\end{equation}
Compactness, weak continuity and finite actions give continuous values and Borel minimizing decisions/actions. Concavity follows by taking infima of the affine risks of fixed continuation programmes: a programme may depend on its subsequently observed reports, but is fixed before the incoming preparation is mixed. Its affine risk need not be continuous for this argument. The instrument identities make that mixture comparison valid. The terminal task and hence $B_n^*$ may change with the horizon.

\begin{definition}[Two finite-state interfaces]\label{def:machines}
An externally clocked observer has finite sets $Z_t$, with $|Z_0|=1$ and $|Z_t|\le m_t$ for $1\le t\le n$. It uses action probabilities $\alpha_t(a\mid z)$ and a Borel stochastic kernel $k_t(z'\mid z,a,y)$. The external phase $t$ is free in this model. Terminal decisions depend only on $Z_n$.

An autonomous observer is a \emph{single} tuple $(Z,z_0,D,\alpha,k,u)$: $|Z|\le M$, $D\subset Z$ is the set of stopping labels, $\alpha(a\mid z)$ is one time-independent action kernel on $Z\setminus D$, $k(z'\mid z,a,y)$ is one time-independent next-label kernel, and $u:D\to U$ is one terminal readout map. At a stopping label it issues the readout and makes no further call. Starting from $z_0$, it must stop after exactly $n$ calls almost surely. Neither kernel may consult absolute time unless it is encoded in $z$. Fresh independent randomization is erased at the next persistent cut. A retained pseudorandom-generator seed, instruction pointer, output tape or clock is part of $Z$.
\end{definition}
A reused autonomous label has one action distribution supported in the intersection of all phase-legal action sets at which it has positive occupancy. The common report space is indexed by action, not by an unread phase. Programme cardinality here bounds labels only; a general Borel rule need not have a finite-bit description. Theorem~\ref{thm:effective} supplies one effective specialization.

A fixed read-only programme can depend on the known preparation, kernels, horizon and budgets, but not on the acquired data or an inaccessible parameter. Programme length, temporary workspace, precision, physical calls and physical time are separate resources. Exact-real measurable programmes define the cardinality theorem; finite-bit implementation is addressed in Section~\ref{sec:resources}. An unused label may have arbitrary programmed behavior. Legality means $\alpha_t(\mathcal A_t\mid z)=1$ on every positive-mass phase label; for a stationary observer the same $\alpha_z$ must obey this requirement at every phase at which $z$ occurs. All measures below are generated by the actual deployed observer, including its fresh randomization.

Write $R_{\ext}(\mathbf m)$ and $R_{\aut}(n,M)$ for the infimal risks in these two classes, with infimum $+\infty$ for an empty class. Write $R_{\ext}(n,M)$ when every $m_t=M$. The checkpoint relaxation permits full history during acquisition and restricts only the final label to $M$ values. These are different information patterns, all using the same preparation and task.

\begin{lemma}[Affine witnesses and executable tests]\label{lem:operational_tests}
On a compact predictive carrier $K$, the uniform closure of the linear span of the constant and the determining affine test coordinates is $\aff(K)$. A finite expression $c_0+\sum_{i=1}^m c_i f_i$ can be realized as the expectation of a signed bounded score by randomly selecting one of the executable tests $f_i$, provided they share a legal continuation budget. The score bound and test-selection cost must be included. A uniform limit is not thereby one finite-cost experiment, and a maximum of affine expectations is not generally an executable single-test expectation.

In the local allocation inequality, uniform approximation of every affine component within $\epsilon$ changes the difference of its two sides by at most $2\epsilon$. In the simultaneous inequality this bound is $2n\epsilon$. Thus every strict finite-coordinate dual violation is preserved by sufficiently accurate finite test-span approximation, without identifying a convex dual function with one physical continuation.
\end{lemma}
\begin{proof}
Let $W$ be the uniform closure of the stated span. If a continuous affine $f$ were not in $W$, Hahn--Banach and the representation of continuous linear functionals on $C(K)$ would give a finite signed measure $\mu$ annihilating $W$ but not $f$. Since constants lie in $W$, its positive and negative parts have equal mass. If that mass is positive, their normalized probabilities have equal determining coordinates and therefore the same barycenter. Affinity gives equal integrals of $f$, a contradiction. Zero mass is immediate. Uniform limits of affine functions are affine, proving equality.

For a finite expression put $C=\sum_i|c_i|$. If $C>0$, choose $i$ with probability $|c_i|/C$, execute its test, and score $c_0+C\operatorname{sgn}(c_i)T_i$, where $T_i$ is the actual test outcome with conditional expectation $f_i$. For normalized $|f_i|\le1$ the score has absolute value at most $|c_0|+C$. The required physical budget is the common allowed budget, with its selector and programme costs, not an assertion that arbitrary collections of differently typed tests have one cheap execution. If $C=0$ only the constant is needed.

The maximum of finitely many functions is 1-Lipschitz in their uniform component norm. The two local measures have mass one. In the simultaneous expression, the left total occupancy mass is $n$, while $\sum_{z,a}\eta_{z,a}(\mathcal Y_a)=n$ and the Radon--Nikodym phase weights sum to one. Applying the same bound on each side gives $2n\epsilon$.
\end{proof}
